\textbf{1. Kiến thức cần nhớ}
\begin{itemize}
\item \textbf{Hằng đẳng thức} là đẳng thức mà hai vế luôn cùng giá trị khi thay các chữ bằng số bất kì.
\item Với $A$, $B$ là các biểu thức tùy ý:
\begin{itemize}
\item Hiệu hai bình phương: $\boxed{A^2-B^2=(A-B)(A+B)}$.
\item Bình phương của một tổng: $\boxed{(A+B)^2=A^2+2AB+B^2}$.
\item Bình phương của một hiệu: $\boxed{(A-B)^2=A^2-2AB+B^2}$.
\end{itemize}
\item Hệ quả hay dùng: $(A-B)^2=(B-A)^2$; \ $(A+B)^2-(A-B)^2=4AB$; \ $(A+B)^2\ge 0$, dấu ``='' khi $A=-B$.
\end{itemize}

\textbf{2. Phương pháp giải các dạng}

%% DANG 8D2?6-1
\textbf{Dạng 1. Khai triển hằng đẳng thức}
\begin{enumerate}
\item Xác định $A$, $B$ (có thể là đơn thức như $2x$, $3y^2$); viết $A$, $B$ trong ngoặc khi có hệ số.
\item Thay vào đúng hằng đẳng thức rồi thu gọn từng hạng tử: $(2x)^2=4x^2$.
\end{enumerate}
Lưu ý: $(A+B)^2\ne A^2+B^2$ — không quên hạng tử $2AB$.

%% DANG 8D2?6-2
\textbf{Dạng 2. Viết biểu thức dưới dạng bình phương, hiệu hai bình phương}
\begin{enumerate}
\item Tìm hai hạng tử là bình phương: $A^2$, $B^2$ (ví dụ $9x^2=(3x)^2$).
\item Kiểm tra hạng tử còn lại có bằng $\pm 2AB$ không; đúng thì viết thành $(A\pm B)^2$.
\item Hai bình phương có dấu trừ ở giữa thì viết $A^2-B^2=(A-B)(A+B)$.
\end{enumerate}
Lưu ý: hạng tử ở giữa phải đúng $2AB$ (không phải $AB$).

%% DANG 8D2?6-3
\textbf{Dạng 3. Tính nhanh, tính giá trị biểu thức}
\begin{enumerate}
\item Tách số thành tổng/hiệu với số tròn: $102^2=(100+2)^2$; \ $49\cdot 51=(50-1)(50+1)$.
\item Áp dụng hằng đẳng thức để tính nhẩm.
\item Tính giá trị biểu thức: rút gọn bằng hằng đẳng thức trước, rồi mới thay số.
\end{enumerate}

%% DANG 8D2?6-4
\textbf{Dạng 4. Chứng minh đẳng thức, tìm giá trị nhỏ nhất, lớn nhất}
\begin{enumerate}
\item Chứng minh đẳng thức: khai triển một vế (thường vế phức tạp) rồi thu gọn bằng vế kia.
\item Tìm GTNN: đưa về dạng $(A\pm B)^2+m\ge m$; GTNN bằng $m$ khi $A\pm B=0$.
\item Tìm GTLN: đưa về dạng $m-(A\pm B)^2\le m$; GTLN bằng $m$ khi $A\pm B=0$.
\end{enumerate}
Lưu ý: luôn chỉ ra giá trị của biến để dấu ``='' xảy ra.
